\documentclass[10pt,a4paper]{amsart}
\usepackage{amsmath,amssymb,amsthm,graphicx}
\usepackage{verbatim}
\newenvironment{reference}{\comment}{\endcomment}
%\newenvironment{reference}{}{}
\newenvironment{slogan}{\comment}{\endcomment}
\newenvironment{history}{\comment}{\endcomment}

% For commutative diagrams we use Xy-pic
\usepackage[all]{xy}

% We use 2cell for 2-commutative diagrams.
\xyoption{2cell}
% Theorem environments.
%
\theoremstyle{plain}
\newtheorem{theorem}[subsection]{Theorem}
\newtheorem{proposition}[subsection]{Proposition}
\newtheorem{lemma}[subsection]{Lemma}

\theoremstyle{definition}
\newtheorem{definition}[subsection]{Definition}
\newtheorem{example}[subsection]{Example}
\newtheorem{exercise}[subsection]{Exercise}
\newtheorem{situation}[subsection]{Situation}

\theoremstyle{remark}
\newtheorem{remark}[subsection]{Remark}
\newtheorem{remarks}[subsection]{Remarks}

\numberwithin{equation}{subsection}

% Macros
%
\def\lim{\mathop{\mathrm{lim}}\nolimits}
\def\colim{\mathop{\mathrm{colim}}\nolimits}
\def\Spec{\mathop{\mathrm{Spec}}}
\def\Hom{\mathop{\mathrm{Hom}}\nolimits}
\def\Ext{\mathop{\mathrm{Ext}}\nolimits}
\def\SheafHom{\mathop{\mathcal{H}\!\mathit{om}}\nolimits}
\def\SheafExt{\mathop{\mathcal{E}\!\mathit{xt}}\nolimits}
\def\Sch{\mathit{Sch}}
\def\Mor{\mathop{\mathrm{Mor}}\nolimits}
\def\Ob{\mathop{\mathrm{Ob}}\nolimits}
\def\Sh{\mathop{\mathit{Sh}}\nolimits}
\def\NL{\mathop{N\!L}\nolimits}
\def\CH{\mathop{\mathrm{CH}}\nolimits}
\def\proetale{{pro\text{-}\acute{e}tale}}
\def\etale{{\acute{e}tale}}
\def\QCoh{\mathit{QCoh}}
\def\Ker{\mathop{\mathrm{Ker}}}
\def\Im{\mathop{\mathrm{Im}}}
\def\Coker{\mathop{\mathrm{Coker}}}
\def\Coim{\mathop{\mathrm{Coim}}}

% Boxtimes
%
\DeclareMathSymbol{\boxtimes}{\mathbin}{AMSa}{"02}

%
% Macros for moduli stacks/spaces
%
\def\QCohstack{\mathcal{QC}\!\mathit{oh}}
\def\Cohstack{\mathcal{C}\!\mathit{oh}}
\def\Spacesstack{\mathcal{S}\!\mathit{paces}}
\def\Quotfunctor{\mathrm{Quot}}
\def\Hilbfunctor{\mathrm{Hilb}}
\def\Curvesstack{\mathcal{C}\!\mathit{urves}}
\def\Polarizedstack{\mathcal{P}\!\mathit{olarized}}
\def\Complexesstack{\mathcal{C}\!\mathit{omplexes}}
% \Pic is the operator that assigns to X its picard group, usage \Pic(X)
% \Picardstack_{X/B} denotes the Picard stack of X over B
% \Picardfunctor_{X/B} denotes the Picard functor of X over B
\def\Pic{\mathop{\mathrm{Pic}}\nolimits}
\def\Picardstack{\mathcal{P}\!\mathit{ic}}
\def\Picardfunctor{\mathrm{Pic}}
\def\Deformationcategory{\mathcal{D}\!\mathit{ef}}
\usepackage{cleveref}
\crefname{originalreference}{original reference}{original references}
\Crefname{originalreference}{Original reference}{Original references}
\makeatletter
\newlabel{definition-split-node}{{0CBV}{source}{}{sourcepointer-1}{}}
\newlabel{definition-split-node@cref}{{[originalreference][0][]0CBV}{[1][0][]source}{}{}{}}
\makeatother
\crefname{originalreference}{original reference}{original references}
\Crefname{originalreference}{Original reference}{Original references}
\makeatletter
\newlabel{lemma-multicross}{{0C1X}{source}{}{sourcepointer-2}{}}
\newlabel{lemma-multicross@cref}{{[originalreference][0][]0C1X}{[1][0][]source}{}{}{}}
\makeatother
\crefname{originalreference}{original reference}{original references}
\Crefname{originalreference}{Original reference}{Original references}
\makeatletter
\newlabel{lemma-split-node}{{0CBW}{source}{}{sourcepointer-3}{}}
\newlabel{lemma-split-node@cref}{{[originalreference][0][]0CBW}{[1][0][]source}{}{}{}}
\makeatother
\title{Derivations of a nodal curve are exactly derivations of its normalization fixing the points over nodes}
\author{The Stacks Project Authors}
\date{stacks-project:curves.tex:local-snapshot-20261012; GNU FDL 1.2 or later}
\begin{document}
\maketitle
\noindent\textit{Editorial note.} All delivered mathematical text is the author's own, copied from the pinned Stacks Project chapter file curves.tex; no statement or proof step was rewritten, reordered or removed.\par\medskip
\noindent\textit{Editorial note.} The counted proof is the author's complete argument: the split-node setup, the node exact sequence, the Leibniz/annihilator computation, the extension fixing x' and x'', and the converse subring preservation are all retained verbatim.\par\medskip
\noindent\textit{Editorial note.} The label lemma-nodal-vector-fields names the counted statement, and the proof environment immediately following it in the source is delivered verbatim as its proof.\par\medskip
\noindent\textit{Editorial note.} The counted statement is carried by the author's own lemma environment in the chapter source.\par\medskip
\noindent\textit{Allowed standard prerequisites (editorial).}\begin{itemize}
\item Split nodes over algebraically closed fields. Same chapter, Definition definition-split-node and Lemma lemma-split-node. Every node of X has two rational inverse images on the normalization since k is algebraically closed.
\item The node exact sequence. Same chapter, Lemma lemma-multicross. Supplies 0 -> O\_\{X,x\} -> O\_\{X\textasciicircum{}nu,x'\} (+) O\_\{X\textasciicircum{}nu,x''\} -> k -> 0, exhibiting the completed local ring at the node inside the product of the two branches.
\end{itemize}\medskip
\ifdefined\hypertarget\hypertarget{sourcepointer-1}{}\fi
\noindent\textit{Source reference.} Original reference 0CBV: sources/stacks-curves/pinned-original-source.tex, source lines 4572--4572. Prerequisite pointer invoked at proof line(s) 6932; it names a standard fact of the same chapter that the counted argument uses, and every step of the counted proof built on it is retained verbatim in the delivered body.\par\medskip
\ifdefined\hypertarget\hypertarget{sourcepointer-2}{}\fi
\noindent\textit{Source reference.} Original reference 0C1X: sources/stacks-curves/pinned-original-source.tex, source lines 3478--3478. Prerequisite pointer invoked at proof line(s) 6940; it names a standard fact of the same chapter that the counted argument uses, and every step of the counted proof built on it is retained verbatim in the delivered body.\par\medskip
\ifdefined\hypertarget\hypertarget{sourcepointer-3}{}\fi
\noindent\textit{Source reference.} Original reference 0CBW: sources/stacks-curves/pinned-original-source.tex, source lines 4585--4585. Prerequisite pointer invoked at proof line(s) 6933; it names a standard fact of the same chapter that the counted argument uses, and every step of the counted proof built on it is retained verbatim in the delivered body.\par\medskip
\begin{lemma}
\label{lemma-nodal-vector-fields}
Let $k$ be an algebraically closed field.
Let $X$ be an at-worst-nodal, proper, connected
$1$-dimensional scheme over $k$. Let $\nu : X^\nu \to X$ be the normalization.
Let $S \subset X^\nu$ be the set of points where $\nu$ is not an
isomorphism. Then
$$
\text{Der}_k(\mathcal{O}_X, \mathcal{O}_X) =
\{D' \in \text{Der}_k(\mathcal{O}_{X^\nu}, \mathcal{O}_{X^\nu}) \mid
D' \text{ fixes every }x^\nu \in S\}
$$
\end{lemma}

\begin{proof}
Let $x \in X$ be a node. Let $x', x'' \in X^\nu$ be the inverse images
of $x$. (Every node is a split node since $k$ is algebriacally closed, see
Definition \ref{definition-split-node} and
Lemma \ref{lemma-split-node}.)
Let $u \in \mathcal{O}_{X^\nu, x'}$ and $v \in \mathcal{O}_{X^\nu, x''}$
be uniformizers. Observe that we have an exact sequence
$$
0 \to \mathcal{O}_{X, x} \to
\mathcal{O}_{X^\nu, x'} \times \mathcal{O}_{X^\nu, x''} \to k \to 0
$$
This follows from Lemma \ref{lemma-multicross}.
Thus we can view $u$ and $v$ as elements of $\mathcal{O}_{X, x}$
with $uv = 0$.

\medskip\noindent
Let $D \in \text{Der}_k(\mathcal{O}_X, \mathcal{O}_X)$. Then
$0 = D(uv) = vD(u) + uD(v)$. Since $(u)$ is annihilator of
$v$ in $\mathcal{O}_{X, x}$ and vice versa, we see that $D(u) \in (u)$ and
$D(v) \in (v)$. As $\mathcal{O}_{X^\nu, x'} = k + (u)$
we conclude that we can extend $D$ to $\mathcal{O}_{X^\nu, x'}$
and moreover the extension fixes $x'$. This produces a $D'$
in the right hand side of the equality. Conversely, given a
$D'$ fixing $x'$ and $x''$ we find that $D'$ preserves
the subring $\mathcal{O}_{X, x} \subset
\mathcal{O}_{X^\nu, x'} \times \mathcal{O}_{X^\nu, x''}$
and this is how we go from right to left in the equality.
\end{proof}

\end{document}
